% Part: first-order-logic
% Chapter: sequent-calculus
% Section: proving-things-quant

\documentclass[../../../include/open-logic-section]{subfiles}

\begin{document}

\olfileid{fol}{seq}{prq}

\olsection{\usetoken{P}{derivation} kanthi Kuantor}

\begin{ex}
Wenehana $\Log{LK}$-!!{derivation} kanggo sekuen $\lexists[x][\lnot !A(x)]
\Sequent \lnot \lforall[x][!A(x)]$.

Nalika nggarap kuantor, kita kudu ngati-ati supaya ora nerak syarat
eigenvariabel, lan kadhangkala iki mbutuhake kita nyoba-nyoba urutan
panrapan inferensi tartamtu. Lumrahe luwih gampang yen aturan kang
kaiket syarat eigenvariabel ditangani luwih dhisik (aturan iku bakal
mapan luwih ngisor ing bukti kang wis rampung). Becike kita uga nyawang
sawetara langkah ing ngarep lan ngira-ira wujude sekuen wiwitan. Ing
kasus iki, wujude kudu kaya $!A(a) \Sequent !A(a)$. Iki ateges nalika
“mbalikke” aturan kuantor, kita kudu milih term kang padha---kang bakal
dijenengi $a$---kanggo aturan $\lforall$ lan $\lexists$. Yen kita
milih term kang beda kanggo saben aturan, pungkasane kita bakal entuk
wujud kaya $!A(a) \Sequent !A(b)$, kang mesthi ora bisa diderivasi.

Kaya lumrahe, kita miwiti kanthi nulis
\begin{prooftree}
\AxiomC{}
\UnaryInf$\lexists[x][\lnot !A(x)] \fCenter \lnot \lforall[x][!A(x)]$
\end{prooftree}
Kita bisa ngetrapake aturan \LeftR{\exists} utawa aturan
\RightR{\lnot}. Amarga aturan \LeftR{\exists} kaiket syarat
eigenvariabel, becike aturan iki ditangani luwih dhisik, mula kita
bakal nindakake iku.
\begin{prooftree}
\AxiomC{}
\UnaryInf$ \lnot !A(a) \fCenter \lnot \lforall[x][!A(x)]$
\RightLabel{\LeftR{\lexists}}
\UnaryInf$ \lexists[x][\lnot !A(x)] \fCenter \lnot \lforall[x][!A(x)]$
\end{prooftree}
Kanthi ngetrapake aturan \LeftR{\lnot} lan \RightR{\lnot} kanthi
mbalikke, kita entuk
\begin{prooftree}
\AxiomC{}
\UnaryInf$\lforall[x][!A(x)] \fCenter !A(a)$
\RightLabel{\LeftR{\lnot}}
\UnaryInf$\lnot !A(a), \lforall[x][!A(x)] \fCenter $
\RightLabel{\LeftR{\Exchange}}
\UnaryInf$\lforall[x][!A(x)], \lnot !A(a) \fCenter $
\RightLabel{\RightR{\lnot}}
\UnaryInf$ \lnot !A(a) \fCenter \lnot \lforall[x] !A(x)$
\RightLabel{\LeftR{\lexists}}
\UnaryInf$ \lexists[x] \lnot !A(x) \fCenter \lnot \lforall[x] !A(x)$
\end{prooftree}
Ing tahap iki, pilihan siji-sijine yaiku ngetrapake aturan
\LeftR{\forall}. Amarga aturan iki ora kaiket watesan eigenvariabel,
kita aman. Elinga, kita ngupaya ngasilake sekuen wiwitan (kanthi wujud
$!A(a) \Sequent !A(a)$), mula nalika aturan ditrapake, kita kudu milih
$a$ minangka argumen kanggo $!A$.
\begin{prooftree}
\Axiom$!A(a) \fCenter !A(a)$
\RightLabel{\LeftR{\lforall}}
\UnaryInf$\lforall[x][!A(x)] \fCenter !A(a)$
\RightLabel{\LeftR{\lnot}}
\UnaryInf$\lnot !A(a), \lforall[x][!A(x)] \fCenter $
\RightLabel{\LeftR{\Exchange}}
\UnaryInf$\lforall[x][!A(x)], \lnot !A(a) \fCenter $
\RightLabel{\RightR{\lnot}}
\UnaryInf$ \lnot !A(a) \fCenter \lnot \lforall[x][!A(x)]$
\RightLabel{\LeftR{\lexists}}
\UnaryInf$ \lexists[x][ \lnot !A(x)] \fCenter \lnot \lforall[x][!A(x)]$
\end{prooftree}
Ing tahap iki wigati, mligine nalika nggarap kuantor, kanggo mriksa
maneh yen syarat eigenvariabel ora diterak. Amarga aturan siji-sijine
kang ditrapake lan kaiket syarat eigenvariabel yaiku
\LeftR{\exists}, lan eigenvariabel~$a$ ora muncul ing sekuen ngisore
(sekuen-pungkasan), mula iki !!{derivation} kang bener.
\end{ex}

\begin{prob}
Wenehana !!{derivation}s kanggo sekuen-sekuen iki:
\begin{enumerate}
\item $\Sequent (\lforall[x][!A(x)] \land \lforall[y][!B(y)]) \lif
\lforall[z][(!A(z) \land !B(z))]$.
\item $\Sequent (\lexists[x][!A(x)] \lor \lexists[y][!B(y)]) \lif
\lexists[z][(!A(z) \lor !B(z))]$.
\item $\lforall[x][(!A(x) \lif !B)] \Sequent \lexists[y][!A(y)] \lif !B$.
\item $\lforall[x][\lnot !A(x)] \Sequent \lnot\lexists[x][!A(x)]$.
\item $\Sequent \lnot\lexists[x][!A(x)] \lif \lforall[x][\lnot !A(x)]$.
\item $\Sequent \lnot\lexists[x][\lforall[y][((!A(x,y) \lif \lnot
!A(y,y)) \land (\lnot !A(y,y) \lif !A(x,y)))]]$.
\end{enumerate}
\end{prob}

\begin{prob}
Wenehana !!{derivation}s kanggo sekuen-sekuen iki:
\begin{enumerate}
\item $\Sequent \lnot\lforall[x][!A(x)] \lif \lexists[x][\lnot!A(x)]$.
\item $(\lforall[x][!A(x)] \lif !B) \Sequent \lexists[y][(!A(y) \lif !B)]$.
\item $\Sequent \lexists[x][(!A(x) \lif \lforall[y][!A(y)])]$.
\end{enumerate}
(Kabeh iki mbutuhake aturan~$\RightR{\Contraction}$.)
\end{prob}

\end{document}
